% ------------------------------------------------------------
% RF Realty CRM — Manual de Usuario
% Compilar con: xelatex (Overleaf usa pdfLaTeX por defecto)
% ------------------------------------------------------------
\documentclass[12pt,oneside]{article}

% ---- Paquetes ----
\usepackage[utf8]{inputenc}
\usepackage[T1]{fontenc}
\usepackage[spanish]{babel}
\usepackage{geometry}
\usepackage{graphicx}
\usepackage{xcolor}
\usepackage{tabularx}
\usepackage{longtable}
\usepackage{enumitem}
\usepackage{hyperref}
\usepackage{fancyhdr}
\usepackage{titlesec}
\usepackage{bookmark}
\usepackage{fontawesome5}

% ---- Configuración de página ----
\geometry{top=2.5cm, bottom=2.5cm, left=2.5cm, right=2.5cm}
\setlength{\parskip}{6pt}
\setlength{\parindent}{0pt}

% ---- Colores ----
\definecolor{primary}{RGB}{37, 99, 235}
\definecolor{secondary}{RGB}{100, 116, 139}
\definecolor{accent}{RGB}{245, 158, 11}
\definecolor{success}{RGB}{16, 185, 129}
\definecolor{danger}{RGB}{239, 68, 68}
\definecolor{bglight}{RGB}{248, 250, 252}
\definecolor{border}{RGB}{226, 232, 240}

% ---- Hipervínculos ----
\hypersetup{
  colorlinks=true,
  linkcolor=primary,
  urlcolor=primary,
  citecolor=primary,
}

% ---- Encabezado y pie ----
\pagestyle{fancy}
\fancyhf{}
\fancyhead[L]{\footnotesize\textcolor{secondary}{RF Realty CRM}}
\fancyhead[R]{\footnotesize\textcolor{secondary}{Manual de Usuario}}
\fancyfoot[C]{\thepage}
\renewcommand{\headrulewidth}{0.5pt}
\renewcommand{\footrulewidth}{0pt}

% ---- Formato de títulos ----
\titleformat{\section}{\Large\bfseries\color{primary}}{}{0em}{}[\vspace{-4pt}\rule{\textwidth}{0.5pt}]
\titleformat{\subsection}{\large\bfseries\color{black!80}}{}{0em}{}
\titleformat{\subsubsection}{\normalsize\bfseries\color{black!70}}{}{0em}{}

% ---- Portada ----
\title{
  \vspace*{3cm}
  {\Huge\bfseries\color{primary} RF Realty CRM}\\[0.5cm]
  {\Large Manual de Usuario}\\[1.5cm]
  {\large Versión 1.0}\\
  {\large Julio 2026}
}
\author{}
\date{}

% ------------------------------------------------------------
\begin{document}

\maketitle
\thispagestyle{empty}
\vfill
\begin{center}
  \textcolor{secondary}{\small Sistema de Gestión Inmobiliaria}\\[4pt]
  \textcolor{secondary}{\small \url{https://real-estate-crm-ten-rosy.vercel.app}}
\end{center}
\newpage

% ---- Índice ----
\tableofcontents
\newpage

% ============================================================
\section{Acerca del Sistema}
% ============================================================

\textbf{RF Realty CRM} es una plataforma web diseñada para que agentes y corredores de bienes raíces puedan gestionar de forma centralizada sus clientes, propiedades, citas y documentos.

\subsection{¿Qué problemas resuelve?}

\begin{tabularx}{\textwidth}{|X|X|}
\hline
\rowcolor{primary!10}
\textbf{Problema} & \textbf{Solución} \\
\hline
Leads perdidos en hojas de cálculo o post-its & Base de datos centralizada con búsqueda instantánea \\
\hline
Seguimiento inconsistente de clientes & Pipeline visual con estados (Nuevo $\rightarrow$ Contactado $\rightarrow$ Calificado $\rightarrow$ Negociación $\rightarrow$ Cerrado) \\
\hline
Propiedades desorganizadas & Catálogo con filtros, tipos, precios y vinculación a portafolios \\
\hline
Citas y llamadas sin registro & Calendario integrado con eventos por tipo y sincronización a Apple Calendar \\
\hline
Documentos dispersos & Repositorio único vinculado a leads, propiedades o portafolios \\
\hline
Atención lenta a leads de WhatsApp & Captura automática con inteligencia artificial y respuesta inmediata \\
\hline
\end{tabularx}

\subsection{Capacidades principales}

\begin{itemize}[leftmargin=*]
  \item \textbf{Gestión de leads:} Alta, edición, eliminación, pipeline Kanban, filtros.
  \item \textbf{Catálogo de propiedades:} CRUD completo, importación desde Excel.
  \item \textbf{Portafolios:} Agrupación de propiedades por criterios (año, tipo, proyecto, etc.).
  \item \textbf{Calendario:} Eventos con tipos (Llamada, Reunión, Visita, Tarea), sincronización WebCal.
  \item \textbf{Documentos:} Subida de archivos PDF, DOC, DOCX, XLS, XLSX con vista previa.
  \item \textbf{WhatsApp AI:} Recepción automática de leads desde WhatsApp con extracción inteligente de datos.
  \item \textbf{Dashboard:} KPIs, gráficos, insights del negocio en tiempo real.
\end{itemize}

% ============================================================
\section{Requisitos y Acceso}
% ============================================================

\subsection{Requisitos técnicos}

\begin{tabularx}{\textwidth}{|X|X|}
\hline
\rowcolor{primary!10}
\textbf{Requisito} & \textbf{Detalle} \\
\hline
Dispositivo & Computadora, tablet o smartphone \\
\hline
Navegador & Google Chrome (recomendado), Mozilla Firefox, Safari, Microsoft Edge \\
\hline
Conexión a internet & Mínimo 1 Mbps (funciona con 3G/4G móvil) \\
\hline
Resolución de pantalla & Mínimo 320px (smartphone), recomendado 1024px o más \\
\hline
Cuenta de usuario & Proporcionada por el administrador del sistema \\
\hline
\end{tabularx}

\subsection{URL de acceso}

\begin{center}
\fbox{\Large\url{https://real-estate-crm-ten-rosy.vercel.app}}
\end{center}

\noindent\textbf{Nota:} El acceso al sistema es privado y por invitación. No es posible registrarse de forma pública. Si necesitas una cuenta, contacta al administrador.

% ============================================================
\section{Inicio de Sesión}
% ============================================================

\subsection{Ingresar al sistema}

\begin{enumerate}[leftmargin=*]
  \item Abre tu navegador web (Chrome, Safari, Edge o Firefox).
  \item En la barra de direcciones, escribe la URL del sistema.
  \item Presiona \textbf{Enter}.
  \item Verás la pantalla de inicio de sesión con los campos \textbf{Email address} y \textbf{Password}.
  \item Ingresa tu correo electrónico.
  \item Ingresa tu contraseña.
  \item Haz clic en el botón \textbf{“Sign in”}.
\end{enumerate}

\subsection{Recuperar contraseña}

Si olvidaste tu contraseña:

\begin{enumerate}[leftmargin=*]
  \item En la pantalla de login, haz clic en \textbf{“Forgot password”} (si está habilitado).
  \item Ingresa tu correo electrónico.
  \item Revisa tu bandeja de entrada: recibirás un enlace para restablecerla.
  \item Si no ves el correo, revisa la carpeta de Spam o Correo no deseado.
\end{enumerate}

\noindent Si la opción “Forgot password” no está disponible, contacta al administrador.

\subsection{Cerrar sesión}

\begin{enumerate}[leftmargin=*]
  \item En la esquina inferior izquierda de la barra lateral, haz clic en tu nombre o avatar.
  \item Selecciona \textbf{“Sign out”}.
  \item Serás redirigido a la pantalla de inicio de sesión.
\end{enumerate}

\subsection{Primera vez que ingresas}

Si es tu primer acceso, el administrador te habrá enviado una invitación por correo electrónico:

\begin{enumerate}[leftmargin=*]
  \item Revisa tu bandeja de entrada: busca un correo de \textbf{Supabase Authentication}.
  \item Haz clic en el enlace \textbf{“Accept Invitation”}.
  \item Establece tu nombre completo y una contraseña segura.
  \item Confirma la contraseña.
  \item Haz clic en \textbf{“Save”} o \textbf{“Confirm”}.
  \item Ve a la URL del sistema e inicia sesión con tu correo y la nueva contraseña.
\end{enumerate}

\noindent\textbf{Recomendación de contraseña segura:} Usa al menos 8 caracteres, combinando mayúsculas, minúsculas, números y símbolos. Ejemplo: \texttt{Casa\#2026!}

% ============================================================
\section{Panel Principal (Dashboard)}
% ============================================================

\subsection{Tarjetas de resumen}

Al iniciar sesión, verás cuatro tarjetas con métricas clave:

\begin{tabularx}{\textwidth}{|X|X|c|}
\hline
\rowcolor{primary!10}
\textbf{Tarjeta} & \textbf{Muestra} & \textbf{Ejemplo} \\
\hline
Total Leads & Número total de clientes potenciales & 24 leads \\
\hline
Properties & Número total de propiedades en el catálogo & 102 propiedades \\
\hline
Portfolios & Número de portafolios creados & 3 portafolios \\
\hline
Events & Eventos del calendario en el mes actual & 12 eventos \\
\hline
\end{tabularx}

\subsection{Gráficos}

\begin{itemize}[leftmargin=*]
  \item \textbf{Gráfico de leads por mes (barras):} Muestra la cantidad de leads creados en cada uno de los últimos 6 meses.
  \item \textbf{Gráfico de propiedades por tipo (anillo):} Distribución porcentual de propiedades según su tipo (House, Apartment, Land, Commercial).
\end{itemize}

\subsection{Insights rápidos}

\begin{tabularx}{\textwidth}{|X|X|}
\hline
\rowcolor{primary!10}
\textbf{Insight} & \textbf{Descripción} \\
\hline
Último lead & Nombre del lead más recientemente creado \\
\hline
Propiedad más cara & Precio y título de la propiedad de mayor valor \\
\hline
Total portafolios & Número total de portafolios activos \\
\hline
\end{tabularx}

% ============================================================
\section{Navegación General}
% ============================================================

\subsection{Barra lateral}

A la izquierda de la pantalla encontrarás la barra de navegación principal:

\begin{tabularx}{\textwidth}{|c|X|X|}
\hline
\rowcolor{primary!10}
\textbf{Ícono} & \textbf{Sección} & \textbf{Descripción} \\
\hline
\faChartBar & Dashboard & Panel principal con resumen del negocio \\
\hline
\faUsers & Leads & Gestión de clientes potenciales \\
\hline
\faHome & Properties & Catálogo de propiedades inmobiliarias \\
\hline
\faFolderOpen & Portfolio & Portafolios de propiedades agrupadas \\
\hline
\faCalendarAlt & Calendar & Calendario de eventos, citas y tareas \\
\hline
\faFileAlt & Documents & Repositorio de documentos \\
\hline
\end{tabularx}

\subsection{Cómo moverse entre secciones}

\begin{enumerate}[leftmargin=*]
  \item Haz clic en cualquier ícono o nombre de sección en la barra lateral.
  \item La sección seleccionada se resaltará con un color diferente.
  \item El contenido principal cambiará automáticamente.
  \item La barra lateral permanece visible para cambiar de sección rápidamente.
\end{enumerate}

\subsection{En dispositivos móviles}

En pantallas pequeñas (smartphones), la barra lateral se oculta automáticamente. Para acceder a ella:

\begin{enumerate}[leftmargin=*]
  \item Toca el ícono de menú (\faBars) en la esquina superior izquierda.
  \item Se desplegará la barra lateral.
  \item Toca la sección a la que deseas ir.
  \item La barra se cerrará automáticamente.
\end{enumerate}

% ============================================================
\section{Gestión de Leads (Clientes Potenciales)}
% ============================================================

\subsection{¿Qué es un lead?}

Un \textbf{lead} es una persona que ha mostrado interés en comprar, vender o alquilar una propiedad. Cada lead contiene: nombre completo, correo electrónico, número de teléfono, estado actual en el proceso de venta, fuente de origen, notas y resumen de IA (si aplica).

\subsection{Vista de lista}

Al entrar en \textbf{Leads}, verás una tabla con todos tus leads:

\begin{tabularx}{\textwidth}{|X|X|X|}
\hline
\rowcolor{primary!10}
\textbf{Columna} & \textbf{Descripción} & \textbf{Ejemplo} \\
\hline
Name & Nombre completo del lead & María García López \\
\hline
Email & Correo electrónico & maria@ejemplo.com \\
\hline
Phone & Teléfono de contacto & +51999000111 \\
\hline
Status & Estado actual (con etiqueta de color) & New \\
\hline
Source & Origen del lead & Website \\
\hline
Created & Fecha de creación & 2026-06-15 \\
\hline
\end{tabularx}

\subsection{Crear un lead nuevo}

\begin{enumerate}[leftmargin=*]
  \item Haz clic en el botón \textbf{“+ Add Lead”} en la parte superior derecha.
  \item Completa el formulario:
\end{enumerate}

\begin{tabularx}{\textwidth}{|X|c|X|X|}
\hline
\rowcolor{primary!10}
\textbf{Campo} & \textbf{Obligatorio} & \textbf{Descripción} & \textbf{Ejemplo} \\
\hline
Name & Sí & Nombre completo del cliente & María García López \\
\hline
Email & Sí & Correo electrónico válido & maria@ejemplo.com \\
\hline
Phone & No & Teléfono con código de país & +51999000111 \\
\hline
Status & Sí & Estado en el pipeline & New \\
\hline
Source & Sí & Cómo llegó el cliente & Website \\
\hline
Notes & No & Notas internas o comentarios & Busca casa en Miraflores, 3 hab \\
\hline
\end{tabularx}

\begin{enumerate}[leftmargin=*, start=3]
  \item Haz clic en el botón \textbf{“Add Lead”}.
  \item El lead aparecerá inmediatamente en la tabla.
\end{enumerate}

\subsection{Ver detalles de un lead}

\begin{enumerate}[leftmargin=*]
  \item En la tabla de leads, haz clic en el \textbf{nombre} del lead.
  \item Se abrirá un panel lateral (drawer) desde la derecha con:
  \begin{itemize}
    \item Información de contacto (email, teléfono, fecha de creación)
    \item Estado actual y fuente de origen
    \item Propiedades de Interés (propiedades vinculadas)
    \item Documentos (archivos asociados)
    \item Notas
    \item Resumen de IA (si aplica)
  \end{itemize}
\end{enumerate}

\subsection{Editar un lead}

\begin{enumerate}[leftmargin=*]
  \item Abre el lead (haz clic en su nombre).
  \item En el panel lateral, haz clic en \textbf{“Edit Lead”}.
  \item Modifica los campos que necesites.
  \item Haz clic en \textbf{“Save Changes”} para guardar o \textbf{“Cancel”} para descartar.
\end{enumerate}

\subsection{Eliminar un lead}

\begin{enumerate}[leftmargin=*]
  \item Abre el lead.
  \item Haz clic en \textbf{“Delete Lead”}.
  \item Aparecerá una ventana de confirmación.
  \item Haz clic en \textbf{“Delete”} para confirmar.
\end{enumerate}

\noindent\textbf{Importante:} La eliminación es permanente. No se puede recuperar un lead eliminado.

\subsection{Pipeline Kanban (vista por etapas)}

El sistema incluye una vista Kanban que organiza los leads en columnas según su estado:

\begin{longtable}{|X|X|X|}
\hline
\rowcolor{primary!10}
\textbf{Columna} & \textbf{Significado} & \textbf{Acción sugerida} \\
\hline
\endhead
New & Lead recién ingresado, sin contacto & Contactar lo antes posible \\
\hline
Contacted & Se ha establecido contacto inicial & Programar seguimiento \\
\hline
Qualified & Cliente calificado con interés real & Enviar propuestas \\
\hline
Negotiation & En negociación activa & Cerrar acuerdo \\
\hline
Closed Won & Venta cerrada exitosamente & Dar seguimiento post-venta \\
\hline
Closed Lost & Venta perdida & Archivar para referencia futura \\
\hline
\end{longtable}

\noindent\textbf{Cómo usar el Kanban:}
\begin{enumerate}[leftmargin=*]
  \item En la página de Leads, haz clic en el botón \textbf{“Kanban”}.
  \item Los leads se organizarán en columnas verticales por estado.
  \item Arrástralos con el mouse de una columna a otra para cambiar su estado.
  \item Para volver a la vista de tabla, haz clic en \textbf{“List”}.
\end{enumerate}

\subsection{Vincular propiedades a un lead}

\begin{enumerate}[leftmargin=*]
  \item Abre el lead (haz clic en su nombre).
  \item En la sección \textbf{“Propiedades de Interés”}, haz clic en \textbf{“+ Agregar”}.
  \item Selecciona la propiedad deseada.
  \item Define el nivel de interés: Bajo, Medio, Alto u Oferta.
  \item Haz clic en \textbf{“Link Property”}.
\end{enumerate}

\noindent\textbf{Para desvincular:} haz clic en el ícono de papelera (\faTrash) junto a la propiedad.

\subsection{Filtrar y buscar leads}

\begin{itemize}[leftmargin=*]
  \item \textbf{Búsqueda por texto:} Escribe en el campo de búsqueda. El sistema buscará en nombre, email y notas.
  \item \textbf{Filtro por estado:} Usa los botones de filtro rápido (All, New, Contacted, etc.).
\end{itemize}

% ============================================================
\section{Gestión de Propiedades}
% ============================================================

\subsection{Lista de propiedades}

En \textbf{Properties} verás el catálogo completo con: título, precio, tipo (House, Apartment, Land, Commercial), estado (Available, Sold, Rented, Under Construction), ciudad, habitaciones, baños y área.

\subsection{Estados de propiedad}

\begin{tabularx}{\textwidth}{|X|X|}
\hline
\rowcolor{primary!10}
\textbf{Estado} & \textbf{Significado} \\
\hline
Available & Disponible para venta o alquiler \\
\hline
Sold & Vendida \\
\hline
Rented & Alquilada \\
\hline
Under Construction & En construcción o proyecto \\
\hline
\end{tabularx}

\subsection{Crear una propiedad}

\begin{enumerate}[leftmargin=*]
  \item Haz clic en \textbf{“+ Add Property”}.
  \item Completa los campos obligatorios: título, precio, tipo.
  \item Opcionalmente agrega: habitaciones, baños, área, dirección, ciudad, descripción, notas.
  \item Haz clic en \textbf{“Add Property”}.
\end{enumerate}

\subsection{Importar propiedades desde Excel}

Puedes cargar múltiples propiedades a la vez usando un archivo \texttt{.xlsx}:

\noindent\textbf{Ejemplo de fila en Excel:}

\begin{verbatim}
| Título               | Precio  | Tipo      | Estado    | Dirección      | Ciudad |
|----------------------|---------|-----------|-----------|----------------|--------|
| Casa en Miraflores   | 250000  | House     | Available | Av. Larco 123  | Lima   |
| Depto en San Isidro  | 180000  | Apartment | Available | Calle Olivos 456 | Lima  |
\end{verbatim}

\noindent\textbf{Pasos:}
\begin{enumerate}[leftmargin=*]
  \item En \textbf{Properties}, haz clic en \textbf{“Import Excel”}.
  \item Selecciona el archivo \texttt{.xlsx}.
  \item El sistema detecta las columnas y crea las propiedades automáticamente.
\end{enumerate}

% ============================================================
\section{Portafolios de Propiedades}
% ============================================================

\subsection{¿Qué es un portafolio?}

Un \textbf{portafolio} es una agrupación de propiedades bajo un mismo tema o categoría. Ejemplos: “Catálogo 2025”, “Casas Premium”, “Inversiones Comerciales”.

\subsection{Crear un portafolio}

\begin{enumerate}[leftmargin=*]
  \item Ve a \textbf{Portfolio}.
  \item Haz clic en \textbf{“+ New Portfolio”}.
  \item Ingresa: nombre, descripción (opcional), año (opcional), tipo, estado.
  \item Haz clic en \textbf{“Create Portfolio”}.
\end{enumerate}

\subsection{Vincular propiedades}

\begin{enumerate}[leftmargin=*]
  \item Abre un portafolio.
  \item Haz clic en \textbf{“+ Add Properties”}.
  \item Selecciona las propiedades que deseas incluir.
  \item Las propiedades aparecerán en la lista del portafolio.
\end{enumerate}

\noindent\textbf{Nota:} Eliminar un portafolio no elimina las propiedades que contiene. Solo se rompe la agrupación.

% ============================================================
\section{Calendario y Eventos}
% ============================================================

\subsection{Vista general}

El calendario muestra una cuadrícula mensual con todos tus eventos programados.

\subsection{Tipos de eventos y colores}

\begin{tabularx}{\textwidth}{|X|X|X|}
\hline
\rowcolor{primary!10}
\textbf{Tipo} & \textbf{¿Para qué usarlo?} & \textbf{Color} \\
\hline
Call & Llamada telefónica con un cliente & Azul \\
\hline
Meeting & Reunión presencial o virtual & Púrpura \\
\hline
Showing & Visita a una propiedad con un cliente & Verde \\
\hline
Task & Tarea pendiente & Ámbar \\
\hline
Other & Cualquier otro tipo de evento & Gris \\
\hline
\end{tabularx}

\subsection{Crear un evento}

\begin{enumerate}[leftmargin=*]
  \item Haz clic en cualquier \textbf{día} del calendario.
  \item Completa: título (obligatorio), hora (opcional), tipo (obligatorio), notas (opcional).
  \item Haz clic en \textbf{“Add Event”}.
\end{enumerate}

\subsection{Eliminar un evento}

Pasa el mouse sobre el evento en el calendario y haz clic en el ícono de papelera (\faTrash) que aparece.

\subsection{Sincronizar con Apple Calendar (WebCal)}

\noindent\textbf{Paso 1: Obtener el enlace}
\begin{enumerate}[leftmargin=*]
  \item En \textbf{Calendar}, haz clic en el botón \textbf{“WebCal”}.
  \item Se copiará un enlace al portapapeles.
\end{enumerate}

\noindent\textbf{Paso 2: En iPhone o iPad}
\begin{enumerate}[leftmargin=*]
  \item Settings $\rightarrow$ Calendar $\rightarrow$ Accounts $\rightarrow$ Add Account $\rightarrow$ Other.
  \item Add Subscribed Calendar.
  \item Pega el enlace y toca Save.
\end{enumerate}

\noindent\textbf{Paso 2: En Mac}
\begin{enumerate}[leftmargin=*]
  \item Calendar $\rightarrow$ File $\rightarrow$ New Calendar Subscription.
  \item Pega el enlace, configura actualización automática, haz clic en OK.
\end{enumerate}

\noindent Los eventos del CRM aparecerán automáticamente en tu calendario de Apple.

% ============================================================
\section{Documentos}
% ============================================================

\subsection{¿Qué son los documentos?}

Los documentos son archivos que puedes subir y asociar a leads, propiedades o portafolios. Cada documento puede ser un PDF, Word o Excel. El sistema permite:

\begin{itemize}[leftmargin=*]
  \item \textbf{Subir} archivos desde tu computadora.
  \item \textbf{Previsualizar} PDFs directamente en el navegador.
  \item \textbf{Descargar} cualquier archivo.
  \item \textbf{Vincular} documentos a leads, propiedades o portafolios.
\end{itemize}

\subsection{Tipos de archivo permitidos}

\begin{tabularx}{\textwidth}{|X|X|c|}
\hline
\rowcolor{primary!10}
\textbf{Tipo} & \textbf{Extensión} & \textbf{Vista previa} \\
\hline
PDF & .pdf & Sí (en navegador) \\
\hline
Word & .doc, .docx & No (solo descarga) \\
\hline
Excel & .xls, .xlsx & No (solo descarga) \\
\hline
\end{tabularx}

\noindent\textbf{Límite:} 10 MB por archivo.

\subsection{Subir un documento}

\textbf{Método 1: Desde la página Documents}
\begin{enumerate}[leftmargin=*]
  \item Ve a \textbf{Documents}.
  \item Haz clic en \textbf{“+ Add Document”}.
  \item Completa: nombre, selecciona archivo, tipo de entidad vinculada (Property, Lead o Portfolio), ID de la entidad.
  \item Haz clic en \textbf{“Upload Document”}.
\end{enumerate}

\textbf{Método 2: Desde el panel de un Lead}
\begin{enumerate}[leftmargin=*]
  \item Abre un lead.
  \item En la sección \textbf{Documents}, haz clic en \textbf{“+ Add”}.
  \item Los campos de entidad ya están prellenados.
  \item Selecciona el archivo y completa los datos.
  \item Haz clic en \textbf{“Upload Document”}.
\end{enumerate}

\subsection{Previsualizar un documento}

\begin{enumerate}[leftmargin=*]
  \item En la lista de documentos (tabla o panel lateral del lead), haz clic en un documento que tenga un archivo subido.
  \item Se abrirá un modal con:
  \begin{itemize}
    \item \textbf{Pantalla completa} para PDFs (se renderiza en el navegador).
    \item \textbf{Botón de descarga} para cualquier tipo de archivo.
    \item \textbf{Botón “Download”} para guardar una copia local.
  \end{itemize}
\end{enumerate}

\subsection{Filtrar documentos}

\begin{itemize}[leftmargin=*]
  \item \textbf{Búsqueda:} Por nombre, descripción o notas.
  \item \textbf{Filtro por entidad:} All Types, Properties, Leads, Portfolios.
\end{itemize}

\subsection{Eliminar un documento}

Haz clic en el ícono de papelera (\faTrash) en la tabla de documentos. El archivo en el almacenamiento también se elimina.

% ============================================================
\section{Integración con WhatsApp}
% ============================================================

\subsection{¿Cómo funciona?}

Cuando un cliente te escribe por WhatsApp, el sistema automatiza la atención:

\begin{enumerate}[leftmargin=*]
  \item \textbf{Recibe} el mensaje de WhatsApp automáticamente.
  \item \textbf{Analiza} el mensaje usando inteligencia artificial (Gemini de Google).
  \item \textbf{Extrae} los datos del cliente: nombre, email, teléfono, presupuesto, tipo de propiedad, ubicación, habitaciones.
  \item \textbf{Crea} un lead en el CRM sin intervención manual.
  \item \textbf{Responde} al cliente confirmando que recibió su información.
\end{enumerate}

\subsection{Ejemplo práctico}

\noindent Mensaje del cliente:
\begin{quote}
“Hola, busco un departamento de 3 habitaciones en San Isidro, mi presupuesto es de 250k. Soy Juan Pérez, mi correo es juan@email.com”
\end{quote}

\noindent Datos extraídos por la IA:

\begin{tabularx}{\textwidth}{|X|X|}
\hline
\rowcolor{primary!10}
\textbf{Campo} & \textbf{Valor extraído} \\
\hline
Nombre & Juan Pérez \\
\hline
Email & juan@email.com \\
\hline
Teléfono & (número de WhatsApp del cliente) \\
\hline
Presupuesto & \$250,000 \\
\hline
Tipo de propiedad & Apartment \\
\hline
Ubicación & San Isidro \\
\hline
Habitaciones & 3 \\
\hline
\end{tabularx}

\subsection{Ver leads de WhatsApp}

Los leads creados por WhatsApp aparecen en \textbf{Leads} con fuente “WhatsApp”. La nota del lead contiene los detalles extraídos.

\subsection{Requisitos técnicos (para el administrador)}

\begin{enumerate}[leftmargin=*]
  \item Obtener una API Key de Gemini en \url{https://aistudio.google.com}.
  \item Configurar Twilio WhatsApp con webhook hacia:
  \begin{center}
  \url{https://real-estate-crm-ten-rosy.vercel.app/api/webhooks/twilio}
  \end{center}
  \item Agregar \texttt{GEMINI\_API\_KEY} en las variables de entorno del hosting.
\end{enumerate}

\subsection{Limitaciones actuales}

\begin{itemize}[leftmargin=*]
  \item La IA funciona mejor con mensajes en español claros y directos.
  \item Si el cliente no proporciona nombre o email, se usan valores genéricos.
  \item La respuesta automática es genérica (no responde preguntas específicas sobre propiedades).
\end{itemize}

% ============================================================
\section{Solución de Problemas}
% ============================================================

\subsection{No puedo iniciar sesión}

\begin{tabularx}{\textwidth}{|X|X|}
\hline
\rowcolor{primary!10}
\textbf{Posible causa} & \textbf{Solución} \\
\hline
Contraseña incorrecta & Usa “Forgot password” para restablecerla \\
\hline
Cuenta no activada & Revisa tu correo (incluyendo Spam) \\
\hline
Error del navegador & Limpia la caché: Configuración $\rightarrow$ Privacidad $\rightarrow$ Borrar datos \\
\hline
\end{tabularx}

\subsection{La importación de Excel falla}

\begin{tabularx}{\textwidth}{|X|X|}
\hline
\rowcolor{primary!10}
\textbf{Problema} & \textbf{Solución} \\
\hline
Archivo no se selecciona & Asegúrate de que sea \texttt{.xlsx} \\
\hline
Columnas no reconocidas & Verifica los encabezados \\
\hline
Precios no se importan & Solo números, sin \$ ni comas \\
\hline
\end{tabularx}

\subsection{El calendario no se sincroniza con iPhone}

\begin{enumerate}[leftmargin=*]
  \item Verifica que el enlace WebCal esté correctamente copiado.
  \item Sigue los pasos exactos de la sección 9.6.
  \item La sincronización puede tomar unos minutos.
  \item Si persiste, elimina la suscripción y vuelve a crearla.
\end{enumerate}

\subsection{No llegan leads de WhatsApp}

\begin{tabularx}{\textwidth}{|X|X|}
\hline
\rowcolor{primary!10}
\textbf{Posible causa} & \textbf{Solución} \\
\hline
Webhook no configurado & El admin debe configurar la URL en Twilio \\
\hline
API Key de Gemini no configurada & Agregar \texttt{GEMINI\_API\_KEY} en variables de entorno \\
\hline
Mensaje muy corto & “Hola”, “Gracias” no generan lead \\
\hline
\end{tabularx}

% ============================================================
\section{Glosario de Términos}
% ============================================================

\begin{longtable}{|X|X|}
\hline
\rowcolor{primary!10}
\textbf{Término} & \textbf{Definición} \\
\hline
\endhead
Lead & Cliente potencial que ha mostrado interés en una propiedad \\
\hline
Pipeline & Proceso de ventas organizado en etapas \\
\hline
Kanban & Método visual donde las tareas se mueven entre columnas por estado \\
\hline
Portafolio & Agrupación de propiedades bajo un mismo tema \\
\hline
WebCal & Estándar de sincronización de calendarios (.ics) \\
\hline
Webhook & Mecanismo que envía datos automáticamente entre sistemas \\
\hline
Gemini & Modelo de IA de Google para analizar mensajes \\
\hline
Twilio & Plataforma de comunicaciones para WhatsApp \\
\hline
KPI & Indicador clave de rendimiento \\
\hline
Entity ID & Identificador único de un lead, propiedad o portafolio \\
\hline
CRUD & Crear, Leer, Actualizar, Eliminar \\
\hline
\end{longtable}

% ============================================================
\section{Soporte y Contacto}
% ============================================================

\subsection{Reportar problemas}

Si encuentras algún error, contacta al administrador proporcionando:
\begin{enumerate}[leftmargin=*]
  \item Descripción del problema: ¿qué estabas haciendo?
  \item Pantallazo (si es posible).
  \item URL exacta de la página.
  \item Navegador y dispositivo.
\end{enumerate}

\subsection{Disponibilidad}

El sistema está alojado en la nube con una disponibilidad objetivo del 99.9\%.

\vspace{1cm}
\begin{center}
\rule{5cm}{0.5pt}\\[4pt]
\textit{Fin del manual}\\[2pt]
RF Realty CRM — Versión 1.0\\[2pt]
Julio 2026
\end{center}

\end{document}
